% ============================================================
%   SESSION 2023
% ============================================================

\sujetbacheader{2023}


% ------------------------------------------------------------
%   EXERCICE 1
% ------------------------------------------------------------

\exercice{3 points}

\noindent\textbf{A.} On considère le système $(S)$ :
$$
\begin{cases}
	3y + 2z = 7 \\
	x + y - z = 2 \\
	2x + y - 3z = 1
\end{cases}
$$

\begin{enumerate}
	\item Écrire $(S)$ sous la forme matricielle : $AX = B$ où $A$ est une matrice carrée d'ordre $3$ et $X$ et $B$ deux matrices colonnes d'ordre $3$. \hfill (0,25 pt)
	\item Soit $P = \begin{pmatrix} -2 & 11 & -5 \\ 1 & -4 & 2 \\ -1 & 6 & -3 \end{pmatrix}$

	Calculer $A \times P$ et $P \times A$. Conclure. \hfill (0,75 pt)
	\item En déduire la matrice $X$. \hfill (0,25 pt)
\end{enumerate}

\bigskip

\noindent\textbf{B.} On dispose deux urnes $U_1$ et $U_2$. L'urne $U_1$ contient trois boules blanches et une boule noire. L'urne $U_2$ contient une boule blanche et une boule noire. Toutes les boules sont indiscernables au toucher. On appelle $E$ l'épreuve suivante : On tire une boule de $U_1$ et on la met dans $U_2$ puis on tire au hasard et simultanément deux boules de $U_2$.

Soit $A$ l'événement : « obtenir deux boules blanches au deuxième tirage. »

\begin{enumerate}
	\item Montrer que la probabilité de l'événement $A$ est égale à $\dfrac{1}{4}$. \hfill (0,75 pt)
	\item On répète trois fois de suite et d'une manière indépendante l'épreuve $E$.

	Soit $X$ la variable aléatoire égale au nombre de réalisations de l'événement $A$.
	\begin{enumerate}
		\item Donner la loi de $X$. \hfill (0,75 pt)
		\item Calculer l'écart-type de $X$. \hfill (0,25 pt)
	\end{enumerate}
\end{enumerate}


% ------------------------------------------------------------
%   EXERCICE 2
% ------------------------------------------------------------

\exercice{2 points}

L'espace est muni d'un repère orthonormé $(O,\vec{\imath},\vec{\jmath},\vec{k})$.

On donne un point $A(2\,;2\,;1)$ et un plan $(\mathscr{P})$ d'équation : $3x + 2y + 6z + 33 = 0$.

\begin{enumerate}
	\item Vérifier que $(\mathscr{P})$ est strictement parallèle à un plan $(Q)$ contenant $A$ et de vecteur normal $\vec{n}(-3\,;-2\,;-6)$. \hfill (0,25 pt)
	\item Déterminer la distance qui sépare ces deux plans. \hfill (0,75 pt)
	\item Soit $H(-1\,;0\,;-5)$.

	Justifier que $H \in (\mathscr{P})$ et que $(AH)$ et $(\mathscr{P})$ sont orthogonaux. \hfill (0,75 pt)
	\item Calculer $\left\Vert \vect{AH} \right\Vert$. \hfill (0,25 pt)
\end{enumerate}


% ------------------------------------------------------------
%   PROBLEME 1
% ------------------------------------------------------------

\probleme{7 points}

Dans le plan orienté $(\mathscr{P})$, on considère un triangle $ABC$, rectangle et isocèle en $A$ tel que $AB = 4\text{cm}$ et $(\vect{AB},\vect{AC}) = \dfrac{\pi}{2}$. On désigne par $O$ le milieu de $[BC]$ et $E$ celui de $[AC]$.

Soient :
\begin{itemize}[label=\textbullet]
	\item $r_1$ : la rotation de centre $A$ et d'angle $\dfrac{\pi}{2}$
	\item $r_2$ : la rotation de centre $O$ et d'angle $\dfrac{\pi}{2}$
\end{itemize}
On pose $f = r_1 \circ r_2$.

\medskip
\noindent\textbf{Partie I}

\begin{enumerate}
	\item
	\begin{enumerate}
		\item Justifier que $f$ est une symétrie centrale. \hfill (0,5 pt)
		\item En décomposant $r_1$ et $r_2$ en deux symétries orthogonales d'axes convenablement choisis, déterminer le centre de $f$. \hfill (0,75 pt)
	\end{enumerate}
	\item Soient $D = bary\{(A,1)\,;(B,-1)\,;(C,-1)\}$ et

	$G = bary\{(B,1-\sqrt{2})\,;(A,\sqrt{2})\}$

	\begin{enumerate}
		\item Déterminer les points $D$ et $G$. \hfill (0,25 pt $\times$ 2)
		\item Déterminer et construire l'ensemble $(\Gamma)$ des points $M$ du plan $(\mathscr{P})$

		vérifiant $MA^2 - MB^2 - MC^2 = -12$. \hfill (0,25 pt $\times$ 2)
	\end{enumerate}
	\item Soit $\overline{S}$ la similitude plane indirecte qui laisse invariant le point $B$ et

	transforme $A$ en $C$. Notons par $F$ le centre de $[CG]$.
	\begin{enumerate}
		\item Vérifier qu'il existe une homothétie $h$ de centre $B$ et transforme

		$A$ en $G$ dont on précisera le rapport. \hfill (0,75 pt)
		\item En déduire les éléments caractéristiques de $\overline{S}$. \hfill (0,75 pt)
	\end{enumerate}
	\item On note par :
	\begin{itemize}[label=--]
		\item $(\mathscr{C})$ le cercle de centre $C$ et de rayon $2$
		\item $I$ le point d'intersection de $(\mathscr{C})$ et $(\Gamma)$ dont $I$ se trouve à l'extérieur du

		carré $ABDC$.
	\end{itemize}
	Montrer que $I$ est le centre de la similitude plane directe $S$, de

	rapport $\sqrt{3}$, d'angle $\dfrac{\pi}{2}$ et qui transforme $C$ en $D$. \hfill (0,75 pt)
\end{enumerate}

\medskip
\noindent\textbf{Partie II}

Le plan complexe est rapporté à un repère orthonormé direct $(A,\vect{AB},\vect{AC})$. On donne
$D(1+\ii)$

\begin{enumerate}
	\item
	\begin{enumerate}
		\item Ecrire l'expression complexe de $r_1$ et $r_2$ \hfill (0,5 pt)
		\item En déduire l'expression complexe, la nature et l'élément caractéristique

		de $f$. \hfill (0,75 pt)
	\end{enumerate}
	\item Déterminer l'expression complexe et les éléments caractéristiques

	de $\overline{S}$. \hfill (0,75 pt)
	\item Déterminer l'expression complexe de $S$ et préciser l'affixe de son

	centre. \hfill (0,5 pt)
\end{enumerate}


% ------------------------------------------------------------
%   PROBLEME 2
% ------------------------------------------------------------

\probleme{8 points}

Pour tout entier $n \in \N^*$, On note par $f_n$ la fonction numérique définie sur $\R$ par
$$
f_n(x) = \e^{nx}\ln(1+\e^{-nx})
$$

On désigne par $(\mathscr{C}_n)$ la courbe représentative de $f_n$ dans un plan muni d'un repère orthonormé $(O,\vec{\imath},\vec{\jmath})$ d'unité $2\text{cm}$.

\medskip
\noindent\textbf{Partie A}

\begin{enumerate}
	\item Calculer les limites de $f_n$ aux bornes de son ensemble de définition

	(Poser $X = \e^{-nx}$) \hfill (0,25 pt $\times$ 2)
	\item On considère la fonction numérique $g_n$ ; $n \in \N^*$ définie sur $\R$ par
	$$
	g_n(x) = \ln(1+\e^{-nx}) - \frac{1}{1+\e^{nx}}
	$$
	\begin{enumerate}
		\item Vérifier que pour tout $x \in \R$, $g_n$ est strictement décroissante \hfill (0,75 pt)
		\item En déduire que pour tout $x \in \R$, $g_n(x) > 0$ \hfill (0,75 pt)
	\end{enumerate}
	\item
	\begin{enumerate}
		\item Montrer que, pour tout $x \in \R$, pour tout $n \in \N^*$ :
		$$
		f_n'(x) = n\e^{nx}\cdot g_n(x)
		$$
		\hfill (0,75 pt)
		\item Dresser le tableau de variation de $f_n$ \hfill (0,75 pt)
	\end{enumerate}
	\item Construire la courbe $(\mathscr{C}_2)$ \hfill (0,75 pt)
	\item En utilisant une intégration par parties, calculer, en $\text{cm}^2$, l'aire du domaine

	plan limité par la courbe $(\mathscr{C}_2)$, l'axe des abscisses et les droites d'équations

	$x = 0$ et $x = 1$. \hfill (0,75 pt)
\end{enumerate}

\medskip
\noindent\textbf{Partie B}

On considère l'équation différentielle $(E) : y' - 2y = \dfrac{-2}{1+\e^{-2x}}$

\begin{enumerate}
	\item Vérifier que l'équation $f_2(x) = \e^{2x}\ln(1+\e^{-2x})$ est une solution de $(E)$

	sur $\R$. \hfill (0,75 pt)
	\item Résoudre dans $\R$, l'équation $(E') : y' - 2y = 0$ \hfill (0,75 pt)
	\item Montrer qu'une fonction $\varphi$ est solution de $(E)$ si et seulement si $(\varphi - f_2)$

	est solution de $(E')$ \hfill (1 pt)
	\item En déduire la solution générale de $(E)$. \hfill (0,5 pt)
\end{enumerate}
